\chapter{Developer guide}
\label{ch:dev}

This chapter is meant for researchers who want to change or extend the code: it
describes how the program is organized, which conventions the routines follow,
and what has to be done to add a time integrator, an analysis type, a load or a
material. \Cref{ch:code} lists the routines, \cref{ch:testing} the tests.

\section{Architecture}

The code has three layers (\cref{fig:flow}):
\begin{enumerate}
  \item \textbf{Driver} (\code{code/Start_Sim.m}, \code{subroutines/Driver}).
    \code{Start_Sim.m} only sets parameters. \code{RunSimulation} reads the
    model, prepares materials and damping, calls the analysis, and saves and
    plots the results. It contains no numerics.
  \item \textbf{Analyses} (\code{subroutines/AnalysisType}). One function per
    analysis type. Each assembles the system, applies loads and boundary
    conditions, solves, and returns a \code{results} struct. They do not plot or
    write files (except the optional solver input of the transient analysis).
  \item \textbf{Building blocks} (all other folders): shape functions,
    quadrature, element matrices, assembly, boundary conditions, readers, time
    integrators, post-processing. They work on plain arrays and know nothing
    about \code{settings}, so they can be reused in studies without the driver.
\end{enumerate}
Data travel only as arguments and outputs. There are no global variables, no
scripts that share the base workspace, and no files that one routine writes and
another reads back. As a result, every function can be called and tested on its
own, and several simulations can run one after another (or in a \code{parfor}
loop) in the same \texttt{MATLAB} session.

\section{Data structures}

\begin{xltabular}{\linewidth}{@{}P{2.4cm}P{3.2cm}L@{}}
  \toprule
  struct & created by & content \\
  \midrule
  \endhead
  \code{settings} & \code{DefaultSettings} & all parameters of a simulation;
    \cref{sec:parameters}. Analyses read it, never change it \\
  \code{model} & \code{ReadModel}, completed by \code{RunSimulation} & mesh of
    the input file and of the spectral elements (\code{Topology_Node},
    \code{xCoor}, \code{yCoor}, \code{p}), sets of loads and boundary conditions
    (with the high-order edge nodes), material data (\code{C}, \code{rho},
    \code{zeta}), Rayleigh parameters, sizes \\
  \code{results} & analysis function & solution and the matrices it was computed
    with; fields in the help text of each analysis function \\
  \code{S} & \code{DynamicAnalysis} & input of \code{Solver_Transient}
    (\code{help}~\code{Solver_Transient}) \\
  \code{p} & \code{settings.p}, \code{model.p} & polynomial degree structure:
    \code{pxi}, \code{peta}, \code{ndist}, \code{metricBasis},
    \code{metricMethod}; passed to every element routine \\
  \bottomrule
\end{xltabular}

\section{Conventions}

\begin{xltabular}{\linewidth}{@{}P{3.6cm}L@{}}
  \toprule
  topic & convention \\
  \midrule
  \endhead
  degrees of freedom & node $n$ has the DOFs $2n-1$ ($u_x$) and $2n$ ($u_y$)
    (\code{Topology_DoF_2D}) \\
  element nodes & corner nodes first, then the edge nodes and the interior nodes
    in the order of \code{ShapeFunctions2D_SEM} (\cref{fig:nodes},
    \code{ReferenceNodes2D}) \\
  input meshes & 4-node or 8-node quadrilaterals; in 8-node meshes the corner
    nodes are numbered before the mid-side nodes \\
  Voigt notation & \texttt{[11 22 33 12 23 13]}, engineering shear strains; plane
    matrices \texttt{[11 22 12]} \\
  units & consistent units of the user's choice (the examples use SI) \\
  matrices & \code{K} and \code{M} are unsymmetric on distorted elements: use
    LU (\texttt{\textbackslash}, \code{decomposition}), never Cholesky \\
  time signals & loads are \code{F + F_dt.Neumann*s(t)}, plus the lifting terms
    of time-dependent prescribed displacements (\code{F_dt.K}, \code{.C},
    \code{.M}) \\
  \bottomrule
\end{xltabular}

\paragraph{Coding style.}
\begin{itemize}
  \item Every routine is a function with a help text: an H1 line
    (\texttt{\%NAME Short description.}) followed by the calling syntax, the
    meaning and size of every input and output, and the method or literature
    reference. \code{CodeHygieneTest} checks that every file has a help text and
    that no two files have the same name (with \code{genpath}, one would shadow
    the other).
  \item New functions use CamelCase names. The names of existing routines are
    kept because studies and the manuscript scripts call them.
  \item Comments explain the reason and the numerical method, not the syntax.
    Keep literature references next to the formula they belong to.
  \item Progress messages go through \code{PrintStep(settings.verbose, ...)} in
    the driver and the analyses; building blocks print little. Plots are made only
    in \code{Visualization} and by \code{RunSimulation}.
  \item The core and the tests use base \texttt{MATLAB} only (the GitHub runner has no
    toolboxes). Toolbox functions are allowed in optional tools (e.g.\
    \code{MMS_2D}); say so in the help text.
  \item \code{.m} files are stored byte for byte (\code{.gitattributes}:
    \code{-text}); many use CRLF line endings and a few are Latin-1 encoded.
    Editors and scripts must preserve both (e.g.\ do not run \texttt{sed -i} on
    them, which converts CRLF to LF).
\end{itemize}

\section{Adding a time integration method}
\label{sec:dev-integrator}

\begin{enumerate}
  \item Write the function with the common calling convention
    (\cref{sec:dynamic}),
\begin{lstlisting}
function [T,U,V,A] = MyMethod(K,M,C,F,F_dt,U0,V0,NumDoF,DOFsave,...
                     ActiveDoFs,ExciteParam,SignalType,dtmax,tsim,myOption)
\end{lstlisting}
    Start from \code{NewmarkMethod} (implicit) or \code{CentralDifferenceMethod}
    (explicit) and use the helpers in \code{TimeIntegration/Common}:
    \code{TimeStepSize} (with the stability factor of an explicit method),
    \code{PrepareSystem}, \code{InitialAcceleration}, \code{SignalHistory},
    \code{LoadAtTime} in the time loop, \code{ExpandToAllDoFs} for the output.
    Only the update formulas of the method remain to be written.
  \item Add one row to the table in \code{TimeIntegrationMethods}: code, name,
    function handle, and a function that returns the trailing arguments from
    \code{settings.TimeIntegration}, e.g.\ \texttt{@(opt) \{opt.Order\}}. Codes
    1xx/2xx are explicit, 3xx/4xx implicit; the first digit also separates
    second-order from higher-order methods.
  \item Add the method to the parameter list of \code{TimeIntegrationTest} with
    its expected order of convergence; the test checks free vibration, a step load
    and the convergence rate on a 2-DOF oscillator.
  \item Add a row to \cref{tab:methods}.
\end{enumerate}

\section{Adding an analysis type}

\begin{enumerate}
  \item Write \code{AnalysisType/MyAnalysis(model, settings)} that returns a
    \code{results} struct. Use \code{AssembleSystem} and \code{LoadVectors} for
    the system and \code{DirichletBCs_homogeneous} for the constrained DOFs, as
    \code{StaticAnalysis} does.
  \item Add its parameters with defaults and comments to \code{DefaultSettings}.
  \item Add a \code{case} to the \code{switch} in \code{RunSimulation} (which
    field to plot, if any) and to \code{SaveResults} (result folder).
  \item Add a test to \code{RunSimulationTest} that compares the result with an
    independent computation.
\end{enumerate}

\section{Adding loads, boundary conditions or materials}

\begin{itemize}
  \item \textbf{Loads.} \code{Read_load_ABAQUS} reads the keyword into
    \code{cload}, \code{dsload}, \code{body_force} or \code{gravity};
    \code{NeumannBCs} loops over them and calls \code{CreateForce},
    \code{CreatePress}, \code{CreateTraction}, \code{CreateBload} or
    \code{CreateGravity}. A new load type needs a reader branch, a
    \code{Create*} routine (edge or area integral with the parametric shape
    functions) and a test in \code{LoadsTest}. Edge loads that are defined on an
    element set get their high-order nodes from \code{pSpectralElement_BCs}.
  \item \textbf{Boundary conditions.} Homogeneous conditions remove DOFs
    (\code{DirichletBCs_homogeneous}); prescribed values are eliminated
    (\code{DirichletBCs_nonhomogeneous}, static) or lifted
    (\code{DirichletBCs_nonhomogeneous_dynamic}, transient).
  \item \textbf{Materials.} \code{Read_material_ABAQUS} stores the constants of
    each material in \texttt{material\_properties\{i,3\}}; \code{MaterialMatrices}
    turns them into the $3\times3$ matrix. A new material model needs a branch in
    both and a test in \code{DriverTest}.
  \item \textbf{Element technology.} The trial space is selected with
    \code{p.metricMethod} in \code{ShapeFunctions2D_Metric}, the basis of its
    Vandermonde matrix with \code{p.metricBasis} (\code{PolynomialBasis1D}), the
    nodes with \code{p.ndist} (\code{nodal_distribution}). New options go into
    these functions and are then available in every analysis; the element tests
    (\code{MetricShapeFunctionsTest}, \code{ElementMatricesTest},
    \code{PatchTest}) can be run for them by adding a parameter value.
\end{itemize}

\section{Using the code in a study}

A study is a folder in \code{studies/} with its own scripts, results and report.
It adds the code to the path (\code{SetupPaths}, or
\code{addpath(genpath('code/subroutines'))}) and then either
\begin{itemize}
  \item loops over \code{RunSimulation} calls with modified settings (complete
    analyses from \texttt{ABAQUS} input files), or
  \item calls the building blocks directly (element-level studies, generated
    meshes), as the existing studies and the helpers in \code{tests/helpers} do.
\end{itemize}
Studies should not modify files in \code{code/}; if they need a change in the
code base, make it there, with a test.

\section{Verifying a change}

\begin{enumerate}
  \item Run \code{run_tests} (about two minutes); GitHub Actions runs it on every
    push.
  \item For changes that should not alter results (refactoring, performance),
    compare the old and the new version bit for bit: run both on the same input
    (e.g.\ the example input files with every analysis type and time integrator)
    and check the outputs with \code{isequal}.
  \item Add a test for every bug that is fixed (\cref{ch:issues}).
\end{enumerate}
